\documentclass[11pt]{article}
\usepackage{fontspec}
\setmainfont{Times New Roman}
\setsansfont{Arial}
\setmonofont{Consolas}
\usepackage{amsmath,amssymb}
\usepackage{geometry}
\usepackage{microtype}
\geometry{a4paper,margin=3cm}
\setlength{\parindent}{1.25em}
\setlength{\parskip}{0pt}
\emergencystretch=30em
\allowdisplaybreaks
\raggedbottom
\providecommand{\p}{\partial}
\providecommand{\frH}{\mathfrak H}
\providecommand{\frK}{\mathfrak K}
\providecommand{\frM}{\mathfrak M}
\providecommand{\frN}{\mathfrak N}
\providecommand{\frG}{\mathfrak G}
\providecommand{\frS}{\mathfrak S}
\providecommand{\frD}{\mathfrak D}
\providecommand{\frf}{\mathfrak f}
\providecommand{\frc}{\mathfrak c}
\providecommand{\fra}{\mathfrak a}
\providecommand{\frb}{\mathfrak b}
\providecommand{\frp}{\mathfrak p}
\providecommand{\frr}{\mathfrak r}
\providecommand{\frm}{\mathfrak m}
\providecommand{\calA}{\mathcal A}
\providecommand{\calB}{\mathcal B}
\providecommand{\calN}{\mathcal N}
\providecommand{\tmod}[1]{\;(\operatorname{mod} #1)}

\begin{document}

% Unit: Paper 14. Direct Ukrainian translation from the German final-audited source slice.

\begin{center}
{\Large\bfseries 14. Арифметична теорія алгебраїчних функцій однієї змінної у її зв'язку з іншими теоріями та з теорією числових полів\par}
\vspace{1.5em}
\emph{Jahresbericht der Deutschen Mathematiker-Vereinigung} 38 (1919), с. 182--203
\end{center}

Наведений нижче огляд виник із бажання мати сполучну ланку між окремими теоріями алгебраїчних функцій; водночас його можна розглядати як доповнення до огляду Брілля--Нетер\footnote{\emph{Die Entwicklung der Theorie der algebraischen Funktionen in älterer und neuerer Zeit}. \emph{Jahresbericht der Deutschen Mathematiker-Vereinigung} т. 3 (1894) [далі цитується як Брілль--Нетер].}, у якому розгляд арифметичної теорії по суті відсутній.\footnote{За винятком обговорення дослідження Кронекера про дискримінант. (\emph{J. f. M.} т. 91.)}

Окрім Брілля--Нетер, для огляду окремих теорій слід мати на увазі такі енциклопедичні статті разом із зазначеною там літературою:

для теорії Рімана: Віртінґер (II B 2), № 1--19;

для теорії Веєрштраса: Віртінґер (II B 2), № 23, 32, 33, 34;

для алгебраїчно-геометричної теорії Брілля--Нетер: Віртінґер (II B 2), № 21--31; Берцоларі (III C 4), № 23--38;

для арифметичної теорії Дедекінда--Вебера і Гензеля--Ландсберга: Гензель (II C 5), [далі цитується як Гензель]\footnote{Частини, що стосуються Дедекінда--Вебера, належать А. Островському.};

загальніше, для арифметичної теорії алгебраїчних елементів: Ландсберг (I C 5), а для арифметичної теорії алгебраїчних функцій кількох змінних: Юнг (II C 6).

Для теорії числових полів пор. «Zahlbericht» Гільберта\footnote{\emph{Die Theorie der algebraischen Zahlkörper}. \emph{Jahresbericht der Deutschen Mathematiker-Vereinigung} т. 4 (1894/95).} і лекції з теорії чисел Діріхле--Дедекінда [далі цитується як Дедекінд--теорія чисел].\footnote{Брауншвейг, Vieweg, 4-е вид., 1894.}

\subsection*{Зміст}
\begin{enumerate}
\item Основи різних теорій.
\item Паралелізм між алгебраїчними числами й алгебраїчними функціями (спряжені елементи, теореми про базис, поняття модуля).
\item Паралелізм і відмінності в теорії ідеалів алгебраїчних чисел і функцій.
\item Зв'язок арифметичного обґрунтування розкладу в ряди з теорією ідеалів.
\item Порівняння арифметичної теорії алгебраїчних функцій з геометричною теорією. Різне поняття інваріантності.
\item Продовження порівняння. Особливі точки.
\item Порівняння між теоремою про залишок і теоремами теорії ідеалів.
\item Трансцендентні питання для алгебраїчних функцій і алгебраїчних чисел.
\end{enumerate}

\subsection*{1. Основи різних теорій}
Теорія Рімана є єдиною суто трансцендентною теорією, що спирається на теореми існування (принцип Діріхле, методи Шварца--Неймана). Інтеграли алгебраїчних функцій тут є первинними, самі алгебраїчні функції --- вторинними. Функції на рімановій поверхні характеризуються своїми нулями та полюсами; це відповідає полігонам (відповідно дивізорам) арифметичної теорії, тоді як Веєрштрас також говорить про самі функції разом з їхніми нулями та полюсами. Головна задача, окрім теорем існування, полягає в побудові всіх функцій із заданими полюсами; її розв'язують за допомогою інтегралів першого і другого роду. Далі постає питання про число довільних сталих, які ще входять при заданих нескінченних точках. На це останнє питання відповідає теорема Рімана--Роха; в арифметичній теорії йому відповідає питання про розмірність класу полігонів (відповідно класу дивізорів), а в алгебраїчно-геометричній теорії --- питання про розмірність повної лінійної системи, на яке також відповідає теорема Рімана--Роха. В арифметичній теорії фактична побудова функцій здійснюється через теореми про базис; в алгебраїчно-геометричній теорії --- через теорему про залишок. Зокрема, нулям підінтегральних функцій першого роду відповідає диференціальний клас в арифметичній теорії, а в алгебраїчно-геометричній теорії --- вчення про спеціальні групи, тобто про систему груп точок, яку вирізають ад'юнктні криві порядку \((n-3)\), \(\varphi\)-криві.

Як ріманова теорія історично передувала іншим теоріям, хоча строге доведення теорем існування було отримане лише пізніше, так і згодом за допомогою її трансцендентних засобів розв'язано алгебраїчні питання, які й досі недоступні суто алгебраїчному або арифметичному трактуванню; тут передусім слід згадати теорію особливих відповідностей Гурвіца.\footnote{A. Hurwitz, Über algebraische Korrespondenzen und das erweiterte Korrespondenzprinzip. \emph{Math. Ann.} т. 28. (Пор. Брілль--Нетер, розд. 10.)}

Аналогії з трансцендентними постановками питань у теорії алгебраїчних числових полів розглянуто в останньому пункті.

Для решти теорій первинними є алгебраїчні функції, а інтеграли --- вторинними. При цьому теорія Веєрштраса спирається на трансцендентні основи розкладу в степеневі ряди і теореми про лишки; далі вона, однак, розвивається алгебраїчно, через явне задання всіх функцій, що розглядаються.

Арифметична теорія у формі Гензеля--Ландсберга також використовує трансцендентні основи розкладу в ряди. Але, зіставляючи на цій основі окремим точкам дивізори, вона далі йде суто арифметичним шляхом і тому може розглядатися як злиття підходів Веєрштраса та Дедекінда--Вебера. На відміну від останніх, вона водночас дає методи фактичної побудови всіх потрібних базисних функцій. У деталях вона має низку спрощень порівняно з Дедекіндом--Вебером, головно завдяки введенню «дробових» дивізорів. Нещодавно Гензель\footnote{\emph{Mathematische Zeitschrift} т. 4; це обґрунтування покладено також в основу його огляду.} арифметично обґрунтував теорію розкладу в ряди (за винятком методу мажорант для доведення збіжності), завдяки чому теорія стала майже суто арифметичною.

Алгебраїчно-геометрична теорія Брілля--Нетер передбачає лише поняття «нескінченно близьких точок» однієї гілки, яке зводиться до розкладу в ряд у звичайній точці\footnote{Для нескінченно близьких «особливих» точок пор. Брілль--Нетер, розд. 6, особливо № 10 і 11. Звичайні нескінченно близькі точки --- це, наприклад, точки перетину дотичної з кривою або точки вищого дотику (точки розгалуження).} і відповідає веєрштрасовому поняттю елемента; далі вона працює алгебраїчно-раціональними методами, по суті методами тернарної теорії елімінування, і також доходить до явних представлень.

Первісна, суто арифметична теорія Дедекінда--Вебера\footnote{\emph{J. f. M.} 92 (цитується як Д.--В.).} споріднена з рімановою остільки, оскільки її твердження знову мають характер доведень існування; вона зберігає повну чистоту методу.\footnote{З цієї причини вона також придатніша для порівняння з іншими теоріями, ніж, скажімо, теорія Гензеля--Ландсберга, яка лише принагідно з'являтиметься в наступних порівняннях. --- Теорія абелевих інтегралів і їх обернення, однак, наскільки вона передбачає неперервність, у Д.--В. і також у подальшому розвитку алгебраїчної теорії (M. Noether, \emph{Ann.} 37) не розглядається.} Арифметичне обґрунтування теорії, де змінні відіграють лише роль невизначених, іде паралельно до теорії алгебраїчних чисел, як докладніше буде пояснено в п. 2.\footnote{Порівняння теорії Гензеля--Ландсберга з теорією чисел, зокрема з теорією \(p\)-адичних чисел, є у Гензеля.} Це обґрунтування дає засіб арифметично ввести поняття точки (пор. Гензель, № 10a) і так побудувати поняття абсолютної ріманової поверхні без жодного міркування про неперервність. Завдяки цьому теорія набуває формального характеру; наприклад, поняття диференціала також вводиться суто формально. На основі, здобутій через поняття точки, можна показати близьку спорідненість із методами й поняттями алгебраїчно-геометричної теорії, яка, якщо не брати до уваги поняття елемента, також вільна від трансцендентних засобів; це буде зроблено в наступних пунктах.

\subsection*{2. Паралелізм між алгебраїчними числами й алгебраїчними функціями. (Спряжені елементи, теореми про базис, поняття модуля)}
Алгебраїчне числове поле означається як поле \(n\)-го степеня над полем \(R\) раціональних чисел; поле алгебраїчних функцій однієї змінної --- як поле \(n\)-го степеня над полем \(K(z)\) усіх раціональних функцій однієї невизначеної \(z\) з довільними комплексними коефіцієнтами. Отже, для обох полів спільно чинні всі теореми, що цілком абстрагуються від спеціальної природи базового поля; це теореми про базис, норми, сліди й дискримінанти.\footnote{Д.--В., § 1 і 2; Гензель, № 4 (з приміткою 5).} У Дедекінда--Вебера всі ці поняття вводяться без використання спряжених елементів. Але слід зауважити, що поняття спряженого поля можна означити цілком формально, як показав Штайніц, спираючись на Кронекера.\footnote{Steinitz, Algebraische Theorie der Körper. (\emph{J. f. M.} 137, § 6 і 8.)} Справді, якщо \(K\) --- довільне поле, а \(\varphi(x)\) --- незвідний над \(K\) поліном від \(x\), то система класів лишків за модулем \(\varphi(x)\) утворює поле. Якщо класові лишків \(x\) поставити у відповідність елемент \(\xi\), то при ізоморфному зіставленні класові лишків, представленому поліномом \(f(x)\), відповідає елемент \(f(\xi)\). Оскільки всі поліноми, кратні \(\varphi(x)\), представляють нульовий клас лишків, \(\varphi(\xi)\) відповідає нульовому класові; тому елемент \(\xi\) можна символічно розглядати як корінь рівняння \(\varphi(x)=0\), і він породжує алгебраїчне поле \(K(\xi)\). Повторення процедури приводить до \(n\) спряжених полів \(K(\xi), K(\xi'),\ldots,K(\xi^{(n-1)})\). Отже, зокрема і для алгебраїчних функцій можна говорити про спряжені елементи, не виходячи за межі суто арифметичної теорії.

Два базові поля \(R\) і \(K(z)\) мають спільну властивість: у них означено систему всіх цілих елементів, відповідно цілих раціональних чисел і цілих раціональних функцій однієї невизначеної. Ці цілі елементи мають характерну властивість: кожні два \(a\) і \(b\) мають найбільший спільний дільник \(d\), який можна подати у формі \(ma+nb\), де \(d,m,n\) також є цілими елементами з \(R\) відповідно \(K(z)\); причому \(d\) є абсолютно найменшим числом відповідно поліномом найнижчого степеня, який можна подати в такій формі. Лише на цій властивості, що тягне єдиність розкладу на прості множники для цілих раціональних чисел відповідно функцій однієї невизначеної, ґрунтуються теореми про цілі елементи алгебраїчного поля або надполя. Міркування, яке доводить існування найбільшого спільного дільника, дає також існування фундаментальної системи, або базису, цілих елементів алгебраїчного поля. Справді, якщо розглядати лише такі базиси поля, які складаються з цілих елементів, то їхній дискримінант є цілим раціональним числом відповідно цілою раціональною функцією однієї невизначеної. Кожен базис, що відповідає абсолютно найменшому числу або функції найнижчого степеня, утворює фундаментальну систему.

Точніше розуміння питань базису, також для загальніших полів, дає введення поняття модуля; його слід сформулювати так, щоб охопити різні означення, які трапляються в літературі залежно від природи елементів. Модулем --- відносно базової області цілісності \(\frS\) --- називатимемо систему елементів таку, що різниця двох елементів знову належить цій системі, і так само добуток елемента на довільний елемент з \(\frS\). (Якщо \(\frS\) складається з цілих раціональних чисел, друга вимога є наслідком першої.)\footnote{Поняття модуля походить від Дедекінда, який уперше ввів числовий модуль (\(\frS=\) цілі раціональні числа) у 2-му виданні «Теорії чисел»; модуль із цілочисельних лінійних форм тут також міститься неявно. У Дедекінда--Вебера з'являються «модулі функцій», елементи яких є алгебраїчними функціями, тоді як \(\frS\) складається з усіх цілих раціональних функцій однієї невизначеної. --- Модуль із однорідних форм від \(n\) невизначених, загальніше з поліномів, уперше з'являється у Гільберта (\emph{Annalen} 36); \(\frS\) складається з усіх поліномів, і, оскільки самі елементи також є поліномами, цей модуль підпадає під наше означення ідеалу. Кронекерові «модульні системи» з поліномів виходять із базису, а не з абстрактного означення.}

Кожен модуль, що має базис із скінченної кількості елементів, через які кожен елемент можна лінійно подати з коефіцієнтами з \(\frS\), називається скінченним (скінченнопородженим); зокрема кожен однопороджений модуль має вигляд \(c\cdot\alpha\), де \(\alpha\) означає базисний елемент, а \(c\) пробігає всі елементи з \(\frS\). \emph{Модуль, елементи якого самі належать \(\frS\), називається ідеалом в \(\frS\);} звідси випливає поняття скінченного ідеалу, зокрема однопородженого (головного ідеалу).

Усі теореми про базис для цілих елементів та ідеалів ґрунтуються на такій теоремі:

\emph{Якщо \(M\) --- модуль відносно \(\frS\), елементи якого є лінійними формами від \(n\) невизначених з коефіцієнтами з \(\frS\), і якщо кожен ідеал у \(\frS\) скінченний, то \(M\) також скінченний модуль. Якщо, зокрема, кожен ідеал у \(\frS\) є головним, то \(M\) має базис із лінійно незалежних форм.}\footnote{У літературі ця теорема з'являється окремо для різних областей коефіцієнтів \(\frS\). Для цілих раціональних чисел пор., наприклад, Гільберт, Zahlbericht § 3 і 4; для поліномів від кількох невизначених: Гільберт, Über die vollen Invariantensysteme: кінець § 2 (\emph{Math. Ann.} 42).}

Справді, якщо елементи \(M\) задані як \(l(x)=a_1x_1+\cdots+a_nx_n\), то сукупність коефіцієнтів \(a_1\) пробігає ідеал у \(\frS\), який за припущенням має вигляд \(b_1a^{(1)}+\cdots+b_\rho a^{(\rho)}\). Належна до \(M\) лінійна форма
\[
m(x)=l(x)-b_1l^{(1)}(x)-\cdots-b_\rho l^{(\rho)}(x)
\]
тоді залежить лише від \((n-1)\) невизначених, звідки скінченним повторенням випливає твердження. Якщо щоразу \(\rho=1\), то базис складається з \(n\) форм відповідно від \(n,n-1,\ldots,1\) невизначених, що дає лінійну незалежність ненульових форм.

Теорема про фундаментальну систему випливає звідси в різних випадках так.

Нехай алгебраїчне поле виникає приєднанням цілого алгебраїчного елемента \(\delta\) до базового поля, і нехай \(\frS\) означає сукупність усіх цілих елементів базового поля. Тоді кожен цілий елемент надполя допускає подання (пор., наприклад, Zahlbericht § 3)
\[
\omega=\frac{a_1\delta^{\,n-1}+a_2\delta^{\,n-2}+\cdots+a_n}{\Delta(\delta)},
\]
де \(a_i\) --- елементи з \(\frS\), а \(\Delta(\delta)\) означає дискримінант \(\delta\). Через підстановку
\[
x_i=\frac{\delta^{\,n-i}}{\Delta(\delta)}
\]
модуль усіх цілих елементів переходить у модуль лінійних форм; отже, існування фундаментальної системи випливає в усіх випадках, коли кожен ідеал у \(\frS\) скінченний. Те саме міркування показує, що загальніше кожен ідеал в області цілих елементів надполя є скінченним. Зокрема для цілих раціональних чисел відповідно функцій однієї невизначеної кожен ідеал є головним; звідси випливає, що фундаментальні системи в алгебраїчних числових полях відповідно в полях алгебраїчних функцій однієї змінної складаються рівно з \(n\) елементів, якщо \(n\) означає степінь. А оскільки в цих полях, за сказаним вище, кожен ідеал скінченний, то водночас випливає фундаментальна система для відносних полів, яка загалом складатиметься з більшої кількості елементів, ніж указує відносний степінь.\footnote{За дослідженнями Штайніца про модулі лінійних форм, де \(\frS\) складається з цілих чисел (скінченного) алгебраїчного числового поля, достатньо \(\nu+1\) базисних елементів, якщо \(\nu\) означає відносний степінь (\emph{Math. Ann.} 71, § 4).} Далі, якщо \(\frS\) означає область усіх поліномів від кількох невизначених, то в \(\frS\) кожен ідеал також є скінченним;\footnote{За гільбертовою теоремою про базис модуля (\emph{Math. Ann.} 36). Модулі з поліномів, які тут заради послідовності названо ідеалами, звичайно називають «модулями форм» або просто «модулями». І гільбертова теорема, що такий модуль \(N\) є скінченним, також зводиться до теореми про лінійні форми. Нехай \(f\) --- поліном \(n\)-го степеня від \(r+1\) невизначених \(x_1,\ldots,x_r,x_{r+1}\), що міститься в \(N\) і має член \(x_{r+1}^{\,n}\) (цього завжди можна досягти лінійним перетворенням). Тоді всі поліноми з \(N\) за модулем \(f\) конгруентні тим, які можна подати у вигляді \(a_1(x)x_{r+1}^{\,n-1}+\cdots+a_n(x)\); отже, через \(x_{r+1}^{\,n-i}=X_i\) вони утворюють модуль лінійних форм, для якого кожен ідеал в області від \(r\) невизначених можна вважати скінченним, бо це відомо для однієї невизначеної.} отже, існування фундаментальної системи випливає і для полів алгебраїчних функцій кількох невизначених; вона знову складатиметься з більшої кількості елементів, ніж указує степінь.

\subsection*{3. Паралелізм і відмінності в теорії ідеалів алгебраїчних чисел і функцій}
Доведення основної теореми теорії ідеалів, що кожен ідеал однозначно подається як добуток степенів простих ідеалів, можна вести спільно для алгебраїчних чисел і функцій, спираючись лише на властивість, що для цілих раціональних чисел відповідно функцій однієї невизначеної кожен ідеал стає головним.\footnote{Пор. зауваження у вступі Гурвіца: Über die Theorie der Ideale (\emph{Gött. Nachr.} 1894). У підручнику алгебри Вебера доведення, у зв'язку з Кронекером і через введення функціоналів, проведено паралельно для обох полів. [Зв'язок між кронекеровою теорією і теорією ідеалів забезпечує «розширена теорема Ґауса»; Гурвіц, там само.]} Головний пункт доведення полягає у встановленні, що з подільності ідеалу \(\frc\) на ідеал \(\fra\) (тобто кожен елемент \(\frc\) міститься в \(\fra\)) випливає добуткове подання \(\frc=\fra\frb\).\footnote{Особливо підкреслено у Дедекінда (Теорія чисел); Supplement XI, Satz VII, § 177; пор. примітку на с. 554. Теорема використовує той факт, що йдеться про алгебраїчні поля; вона вже не чинна для ідеалів (модулів) з поліномів, де на місце добутку стає найменше спільне кратне.} У Гурвіца це доведення спирається на «розширену теорему Ґауса», яка стверджує: якщо цілий елемент \(\omega\) ділить кожен коефіцієнт добутку двох поліномів \(\varphi\) і \(\psi\), то він ділить також добуток усіх коефіцієнтів \(\varphi\) і \(\psi\); у Дедекінда --- на модульну теорему, еквівалентну цій.\footnote{Теорія чисел, Suppl. XI; Satz VI, § 173. На еквівалентність обох теорем указує Дедекінд (Über die Begründung der Idealtheorie, \emph{Göttinger Nachr.} 1895).}

Далі спільно можна означити поняття комплементарного модуля (Гензель, № 12) і з нього --- ідеал розгалуження (основний ідеал, диференту) \(\frD\) (норма якого стає дискримінантом поля \(D\)); і чинна основна теорема, що ідеал розгалуження \(\frD\) є \emph{найбільшим спільним дільником усіх головних ідеалів \(F'(\delta)\)}, де \(\delta\) пробігає всі цілі елементи поля, а \(F(\delta)=0\) кожного разу означає відповідне незвідне рівняння. Звідси випливає, що \emph{дискримінант поля містить усі й лише ті раціональні прості числа відповідно лінійні множники за \(z\), які діляться на квадрат простого ідеалу}. Для самого \(F'(\delta)\) маємо подання
\[
F'(\delta)=\frf\cdot\frD,
\]
і після взяття норм
\[
\Delta(\delta)=R^2\cdot D;
\]
тут \(\frf\) означає провідник кільця (порядку, області цілісності), виведеного з \(\delta\), тобто \(\frf\) складається із сукупності (цілих) елементів, які після множення на кожен цілий елемент діляться на модуль \([1,\delta,\ldots,\delta^{n-1}]\). \(R^2=\operatorname{Norm}\frf\) є квадратом детермінанта підстановки, що переводить базис \(1,\delta,\ldots,\delta^{n-1}\) у базис цілих елементів.\footnote{Dedekind, Über die Diskriminante endlicher Körper (\emph{Abhandl. der Gött. Ges. d. Wissensch.} т. 29, 1882). Подані там для алгебраїчних чисел означення безпосередньо переносяться на алгебраїчні функції; доведення також переносяться майже дослівно. Порядок доведень у Д.--В. від цього відрізняється: див. нижче. Щодо наведених теорем пор. також Zahlbericht, § 31 і 32. (\(F'(\delta)\) там названо «диферентою» цілого числа \(\delta\).)}

\emph{Відмінності} в подальшому розвитку зумовлені спеціальними властивостями базових полів. Так, у випадку алгебраїчних чисел той факт, що \emph{елементи базового поля водночас є показниками}, призводить до того, що ідеал розгалуження може ділитися на вищий, ніж \((e-1)\)-й, степінь простого ідеалу \(\frp\), який входить у \(p\) у \(e\)-му степені (саме коли \(e\) ділиться на \(p\)); це веде до дальших понять поля інерції та поля розгалуження,\footnote{Пор. Zahlbericht, § 39--45.} які не мають аналога в теорії алгебраїчних функцій. Передусім же через \emph{існування нескінченно багатьох сталих} у базовому полі для алгебраїчних функцій зникає скінченність числа класів ідеалів, а разом із нею всі питання, пов'язані з визначенням числа класів. У певному сенсі, однак, як трансцендентний аналог поля класів можна розглядати область усіх (трансцендентних) прайм-форм\footnote{Klein, Zur Theorie der Abelschen Funktionen (\emph{Ann.} 36, 1890).}, бо тут кожна точка породжується однією-єдиною формою, головним ідеалом.

З іншого боку, існування нескінченно багатьох сталих у базовому полі \(K(z)\) для алгебраїчних функцій веде до теорем і постановок питань, що не мають аналога для алгебраїчних чисел. Спершу звідси випливає певне спрощення у побудові доведень; наприклад, у доведенні Д.--В. єдиності розкладу ідеалів на прості ідеали можна оминути «розширену теорему Ґауса» або еквівалентну їй теорему, використовуючи той факт, що кожна лінійна форма \((z-c)\) має нескінченно багато класів лишків, а саме всі сталі (тоді як число класів лишків за модулем простого числа скінченне).\footnote{Д.--В., § 9. Пор. примітки на с. 212 і 213.}

До справді нових результатів веде дальший факт: кожен поліном із коефіцієнтами, що є цілими раціональними функціями від \(z\), за модулем \((z-\alpha)\) розкладається на лінійні множники (чого не буває для цілочисельного полінома за модулем \(p\)); цей факт зберігається навіть тоді, коли замість усіх сталих беруть лише всі алгебраїчні числа.\footnote{Д.--В. у вступі зауважують, що в цьому випадку вся теорія зберігається.} Звідси випливає, що кожен простий ідеал є ідеалом першого степеня;\footnote{Д.--В., § 9, № 7.} а далі --- що дискримінант поля \(D\) стає перетином усіх дискримінантів \(\Delta(\delta)\), де \(\delta\) пробігає всі цілі елементи;\footnote{Д.--В., с. 224.} обидва твердження контрастують із випадком алгебраїчних чисел. Після цього знову випливає, що ідеал розгалуження є перетином усіх головних ідеалів \(F'(\delta)\).

Подальші відмінності зумовлені тим, що для алгебраїчних чисел базове поле \(R\) раціональних чисел є абсолютно виділеним, а саме простим полем, що міститься в даному полі (тобто полем, виведеним з одиниці); натомість базове поле \(K(z)\) є відносним. Кожну несталу функцію \(\xi\) можна вибрати як незалежну змінну; тоді поле стає алгебраїчним полем над \(K(\xi)\). Теорема, що кожен простий ідеал \(\frp\) є ідеалом першого степеня, тобто кожна функція за модулем \(\frp\) конгруентна сталій, разом із можливістю взяти довільну функцію за незалежну змінну веде до найважливішого поняття, що виходить за межі алгебраїчних чисел: до поняття точки в суто арифметичній формі. (Пор. Гензель, 10a.) «Кожна точка тут представлена полем сталих, ізоморфним функціональному полю.»\footnote{Два поля, як відомо, називаються «ізоморфними», якщо їх можна взаємно однозначно зіставити так, що також відповідають сума, різниця, добуток і частка.} Сукупність означених так точок утворює «абсолютну ріманову поверхню». Це означення точки залежить лише від поля, а не від спеціальної змінної. Натомість доведення існування проводиться з виділенням якоїсь функції як незалежної змінної \(z\); точка \(P\) породжується простим ідеалом \(\frp\), коли кожній функції зіставляється її сталий лишок за модулем \(\frp\). Отже, зокрема, функції, подільні на \(\frp\), зникають у \(P\), і навпаки. Коли \(\frp\) пробігає всі прості ідеали за \(z\), цим отримуються всі точки, в яких \(z\) скінченна; решта точок виникає через введення незалежної змінної \(\xi=1/z\) і відповідає простим ідеалам, що входять у головний ідеал \(\xi\). Природно, поняття, пов'язані з поняттям точки, також не можуть мати аналога для алгебраїчних чисел: наприклад поняття полігона, класу полігонів, розмірності класу полігонів тощо.

Слід ще зауважити, що вільний вибір незалежної змінної дає ще одне тлумачення розкладу \(F'(\vartheta)=\frf\cdot\frD\). Якщо розглядати \(\vartheta\) як незалежну змінну, то відповідно маємо
\[
\left(\frac{\p F}{\p z}\right)=\frf\cdot\fra,
\]
оскільки кільце, виведене з \(\vartheta\), було означене симетрично щодо \(z\) і \(\vartheta\); і \(\frf\) стає найбільшим спільним дільником двох головних ідеалів \(\frac{\p F}{\p z}\) та \(\frac{\p F}{\p\vartheta}\). Провідник \(\frf\) водночас стає «ідеалом подвійних точок» щодо \(\vartheta,z\).\footnote{Д.--В., с. 263.}

\subsection*{4. Зв'язок арифметичного обґрунтування розкладу в ряди з теорією ідеалів}
У Гензеля--Ландсберга на місце теорії ідеалів стає трансцендентний допоміжний засіб розкладу в ряди та запозичене з теорії функцій поняття точки. Нещодавнє арифметичне обґрунтування розкладу в ряди для алгебраїчних чисел і функцій, дане Гензелем,\footnote{\emph{Eine neue Theorie der algebraischen Zahlen} (\emph{Math. Zeitschr.} т. 2, 1918) і \emph{Neue Begründung der arithmetischen Theorie der algebraischen Funktionen einer Variablen} (\emph{Math. Zeitschr.} т. 4, 1919). Для докладнішого викладу пор. Гензель, № 44 і далі.} можна тому розглядати як еквівалент теорії ідеалів. Фактично це арифметичне обґрунтування зводиться до того, щоб для кожного простого ідеалу перейти до такого (трансцендентного) поля-розширення, в якому цей простий ідеал стає головним, так що чинні звичайні теореми розкладу; достатньо провести це міркування для всіх простих ідеалів, що входять у просте число \(p\) відповідно в лінійну функцію \(z-a\).

Спершу йдеться про введення трансцендентного поля-розширення, утвореного всіма розкладами в ряди за цілими степенями \(z-a\) відповідно \(p\) (\(p\)-адичними числами), що мають щоразу не більш ніж скінченну кількість від'ємних степенів. У цьому полі незвідне рівняння, яке означає алгебраїчне числове або функціональне поле, розкладається на добуток \(\rho\) незвідних множників; це відповідає розкладу \(p\) відповідно \(z-a\) на добуток \(\rho\) «первинних» ідеалів \(\mathfrak q_1,\ldots,\mathfrak q_\rho\). Дальшому поданню кожного первинного ідеалу як степеня простого ідеалу
\[
\mathfrak q_i=\frp_i^{\,\ell_i},
\]
у Гензеля відповідає введення ще одного поля, алгебраїчного над трансцендентним полем-розширенням. У випадку алгебраїчних функцій його можна означити чистим рівнянням; у випадку алгебраїчних чисел --- загальніше алгебраїчно розв'язним рівнянням. Спрощення для алгебраїчних функцій пояснюється тим, що тут не з'являються поля інерції та розгалуження (пор. п. 3).

«Обривний ряд» за простим елементом \(\pi\),\footnote{\emph{Neue Begründung \(\ldots\)}, формула (11).} що виникає в цих дослідженнях (і ще не потребує трансцендентного доведення збіжності),
\[
v=c_0+c_1\pi+\cdots+c_{\sigma-1}\pi^{\sigma-1}-v_\sigma\pi^\sigma,
\]
де \(c_0,\ldots,c_{\sigma-1}\) означають сталі, а \(v_\sigma\) --- цілий елемент поля, має цілком відповідник у Д.--В.\footnote{Д.--В., с. 231 і § 15.} Там він виводиться з теорії ідеалів і дає основу для означення порядку зникання (відповідно прямування до нескінченності) функції в точці.

\subsection*{5. Порівняння арифметичної теорії алгебраїчних функцій з геометричною теорією. Різне поняття інваріантності}
«Інваріантним» є кожне поняття, охарактеризоване самим лише полем. В арифметичній теорії ця інваріантність загалом уже дана в означенні, яке формулюється щодо всіх елементів поля, а не щодо якогось базису чи виділеної змінної; прикладом є наведене вище поняття точки. Натомість інші теорії виходять із якоїсь спеціальної змінної або базису і потім показують незалежність поняття від цього спеціального вибору; інваріантність тут стає інваріантністю щодо біраціонального перетворення.

Нехай, наприклад, поле з одного боку виникає приєднанням алгебраїчного елемента \(s\) до \(K(z)\), а з іншого --- приєднанням \(\sigma\) до \(K(\xi)\), де \(f(z,s)=0\) відповідно \(F(\xi,\sigma)=0\) означають незвідні рівняння для \(s\) відповідно \(\sigma\); або ще загальніше --- приєднанням скінченної кількості елементів до \(K(\eta)\), з відповідними незвідними рівняннями
\[
g_1(\eta,\tau_1)=0,\qquad g_2(\eta,\tau_1,\tau_2)=0,\ldots .
\]
Тоді, за означенням, усі елементи раціонально виражаються через \(z\) і \(s\), так само через \(\xi\) і \(\sigma\), а також через \(\eta,\tau_1,\tau_2,\ldots\). Зокрема, \(z\) і \(s\) раціонально виражаються через \(\xi\) і \(\sigma\), і навпаки; так само \(z\) і \(s\) раціонально виражаються через \(\eta,\tau_1,\tau_2,\ldots\), і навпаки. Отже, плоска крива \(f(z,s)=0\) переходить через «біраціональне перетворення» у плоску криву \(F(\xi,\sigma)=0\) або також у криву в просторі змінних \(\eta,\tau_1,\tau_2,\ldots\), означену рівняннями \(g_1(\eta,\tau_1)=0; g_2(\eta,\tau_1,\tau_2)=0,\ldots\). Інваріантність виявляється так: означення, сформульоване для спеціальної кривої, зберігається для всіх кривих, що виникають із неї біраціональним перетворенням (у просторі довільної кількості вимірів), тобто для всіх кривих «класу» за Ріманом. «Точка» тут означається як узгоджена система значень \(z=a,s=b\), так що \(f(a,b)=0\); ця точка при біраціональному перетворенні загалом знову переходить в узгоджену систему значень \(\alpha,\beta\) на \(F(\xi,\sigma)=0\), але в особливих випадках, коли \(a,b\) була «особливою» точкою, вона може відповідати кільком «точкам» на \(F(\xi,\sigma)=0\). Завжди можна вказати криві, на яких «особлива» точка \(f(z,s)=0\) відповідає скінченній кількості простих точок, розділених або сусідніх (розв'язання особливостей).\footnote{Пор., наприклад, Брілль--Нетер, розд. 6. Докладніше про «особливі» точки у наступному пункті.} Так само група точок переходить знову в групу точок; число її точок є інваріантним, якщо особливу точку рахувати стількома точками, скільком простим точкам вона відповідає при придатному перетворенні. Оскільки кожну криву можна біраціонально перетворити на таку, що має «звичайні кратні» точки,\footnotemark[\value{footnote}] геометрична теорія розвиває свої поняття лише для таких спеціальних кривих і потім показує інваріантність щодо кожного біраціонального перетворення, також такого, що веде до найзагальніших особливих точок.

Зокрема, як виділену криву «класу» --- за винятком гіпереліптичного випадку --- можна ввести так звану «нормальну криву \(\varphi\)» у просторі \((p-1)\) вимірів, яка не має особливих точок; тут \(\varphi\) --- це \(p\) «ад'юнктних кривих порядку \((n-3)\)» для однорідненої \(f(z,s)=0\). Біраціональному перетворенню \(f(z,s)=0\) у \(F(\xi,\sigma)=0\) відповідає лінійне перетворення \(\varphi\), так що тут ідеться про інваріантність у проєктивному сенсі. Подання всіх функцій поля як раціональних функцій від \(\varphi\) у геометричній теорії називається «інваріантним поданням».\footnote{Пор., наприклад, Брілль--Нетер, розд. 8. Для алгебраїчних функцій кількох змінних питання, чи можна інваріантність звести до інваріантності щодо лінійного перетворення, не розглядається.}

Отже, лінійна система \(\varphi\), або також вирізані ними «спеціальні групи», при кожному біраціональному перетворенні переходить у себе, а тому характеризується самим лише полем. Вона відповідає диференціальному класу в арифметичній теорії, де інваріантність знову безпосередньо виявляється в означенні, даному через властивості щодо кожної функції поля як незалежної змінної.

\subsection*{6. Продовження порівняння. Особливі точки}
З різним розумінням поняття точки пов'язано те, що «особливі точки» кривих, які в алгебраїчно-геометричній теорії створюють особливі труднощі, в обґрунтуванні арифметичної теорії взагалі не з'являються; і тому «ад'юнктні функції» також не потрібні для побудови арифметичної теорії, як це має місце в геометричній теорії, а належать до спеціальних вищих частин теорії. З іншого боку, у фундаменті геометричної теорії не з'являються точки розгалуження; фактично вони зникають і в поданні інтегралів в аронгольдовій нормальній формі (Гензель, № 27, формула 80).

Означення особливої точки, дане в п. 5, можна сформулювати так: \(f(z,s)=0\) (відповідно крива у вищому просторі) має в \(z=a,s=b\) (відповідно \(z=a; s_i=b_i\)) особливу точку, якщо цю систему значень приймає більш ніж одна точка абсолютної ріманової поверхні (точка в арифметичному сенсі), або якщо вона приймається в одній чи кількох точках у вищому порядку. В останньому випадку йдеться про точку полігона розгалуження як для \(z\), так і для \(s\); це відповідає «нескінченно близьким точкам» геометричної теорії; точка стає (вищою) точкою повернення кривої \(f(z,s)=0\) (відповідно кривої у вищому просторі).

Оскільки кожна довільна функція \(\eta\) поля за припущенням раціонально виражається через \(z\) і \(s\), кожна точка абсолютної ріманової поверхні, де \(z=a,s=b\), через ізоморфізм виникає так: кожній функції \(\eta\), поданій через \(z\) і \(s\), зіставляється її значення при \(z=a,s=b\). Щоб точка була особливою, має існувати принаймні одна функція \(\eta\), яка при \(z=a,s=b\) набуває кількох значень або кілька разів тієї самої системи значень. А через раціональну виражуваність через \(z\) і \(s\) це можливо лише тоді, коли в цьому поданні при \(z=a,s=b\) чисельник і знаменник одночасно зникають. Звідси випливає, що наведене означення тотожне звичайному, за яким
\[
f,\quad \frac{\p f}{\p z},\quad \frac{\p f}{\p s}
\]
зникають при \(z=a,s=b\). Адже в цьому випадку
\[
\eta=\frac{\frac{\p f}{\p z}}{\frac{\p f}{\p s}}
\]
є функцією зазначеного типу, багатозначною при \(z=a,s=b\), тоді як у протилежному випадку кожна функція, навіть якщо вона стає \(0/0\), набуває лише одного значення.

За арифметичного поняття точки особлива точка не може з'явитися, бо дві точки вважаються різними, щойно бодай одній функції відповідають різні системи значень. Але й у теорії ідеалів, що служить для обґрунтування арифметичного поняття точки, поняття особливої точки не виникає, оскільки завжди одночасно розглядається сукупність усіх цілих елементів поля. Усі особливі точки \(F(z,\delta)=0\), що лежать у скінченній частині, де \(\delta\) є алгебраїчно цілою щодо \(z\), породжуються (за другим означенням і кінцем п. 3) «ідеалом подвійних точок» \(\frf\). Оскільки за основною теоремою ідеал розгалуження \(\frD\) є найбільшим спільним дільником усіх головних ідеалів \(F'(\delta)\), то для кожного простого ідеалу \(\frp\) можна вказати таке \(\delta\), що відповідний \(\frf\) не ділиться на \(\frp\); отже, оскільки \(\frf\) є найбільшим спільним дільником \(F'(\delta)\) і \(F'(z)\), принаймні один із двох головних ідеалів \(F'(\delta)\) та \(F'(z)\) не ділиться на \(\frp\). Тому \emph{не існує системи значень, для якої одночасно зникають усі \(F'(\delta)\) і \(F'(z)\)}, і, отже, система значень \(z=a,\delta_i=b_i\) приймається лише в одній точці абсолютної ріманової поверхні; \emph{сукупність цілих алгебраїчних функцій поля} (яку можна тлумачити як криву в просторі нескінченно багатьох вимірів) \emph{не має особливої точки у скінченній частині}; у теорії ідеалів же розглядаються лише скінченні значення.

Базисні функції \(\omega_1,\ldots,\omega_n\) усіх цілих елементів також не мають особливих точок у скінченній частині. Справді, за сказаним вище, кожна точка в арифметичному сенсі вже однозначно визначена системою сталих, ізоморфно зіставленою всім цілим функціям. Якщо, отже, «просторова крива», означена якимись цілими функціями \(s_i\), має особливу точку у скінченній частині при \(s_i=\alpha_i\), то має існувати принаймні одна ціла функція \(\eta\), яка при \(s_i=\alpha_i\) набуває кількох систем значень. Якщо тепер \(\omega_1,\ldots,\omega_n\) означає базис усіх цілих елементів, то через \(\delta=a_1\omega_1+\cdots+a_n\omega_n\) з цілими раціональними \(a_i\) кожній системі значень \(\omega\) відповідає лише одна система значень \(\delta\); отже, просторова крива, означена \(\omega_1,\ldots,\omega_n\), не може мати особливої точки у скінченній частині. І \emph{кожна точка абсолютної ріманової поверхні, в якій \(z\) скінченна, вже однозначно визначена системою сталих, ізоморфно зіставленою \(\omega\). Базисні функції --- це саме ті функції \(\eta\), які стають багатозначними при \(s_i=\alpha_i\).} Отже, введення базису цілих елементів, наприклад замість базису \(1,\delta,\ldots,\delta^{n-1}\), обходить труднощі особливих точок.

У геометричній теорії задача однозначного породження всіх точок абсолютної ріманової поверхні --- загальніше, побудови всіх функцій, які стають нескінченними в заданих точках, --- виконується за допомогою ад'юнктних форм та їхніх часток. Форму \(\psi\) називають ад'юнктною, якщо \(\psi=0\) має в кожній \(i\)-кратній точці базової кривої \(f(z,s)=0\) принаймні \((i-1)\)-кратну точку; вищу особливу точку при цьому спершу треба розкласти на звичайні кратні точки. Вище було показано, що для найзагальнішої функції \(\eta=\frac{\psi}{\chi}\) у кожній особливій точці \(f(z,s)=0\) чисельник і знаменник взагалі мусять зникати; те, що вони мусять бути ад'юнктними, показує розгляд усюди скінченних диференціалів. Обернено, що умови ад'юнктності також достатні для побудови найзагальнішої функції, показує «теорема про залишок», до якої ми ще докладніше повернемося. Отже, ад'юнктні функції заступають місце базисних функцій арифметичної теорії.

Арифметично ад'юнктній формі відповідає чисельник функції поля, подільної на «ідеал подвійних точок» \(\frf\), щойно особливі точки --- чого завжди можна досягти перетворенням --- вважаються всі розташованими у скінченній частині. Звідси формально видно подання базисних елементів \(\omega_1,\ldots,\omega_n\) як часток ад'юнктних форм, а отже і зв'язок між різними способами їх подання:

Нехай \(R(z)\) означає детермінант підстановки базису \(1,\delta,\ldots,\delta^{n-1}\) у базис \(\omega_1,\ldots,\omega_n\). Тоді, навпаки, кожне \(\omega\) виражається як частка цілої функції від \(z\) і \(\delta\) на \(R(z)\):
\[
\omega_i=\frac{G_i(z,\delta)}{R(z)}.
\]
Оскільки \(R(z)\omega_i\) належить кільцю, виведеному з \(\delta\), то за означенням провідника \(R(z)\) ділиться на \(\frf\) [із \((R(z))^2=\operatorname{Norm}\frf\) ця подільність випливає лише для \((R(z))^2\)]. А оскільки \(\omega_i\), як ціла функція, лишається скінченною для всіх скінченних \(z\), то і чисельникова функція \(G_i(z,\delta)\) мусить ділитися на \(\frf\); отже, чисельники й знаменники всіх \(\omega_i\) є ад'юнктними. З цього подання \(\omega\) випливає подання всіх функцій поля частками ад'юнктних.

З огляду на попереднє геометричну теорію можна охарактеризувати як \emph{теорію кільця, виведеного з елемента \(\delta\)} (порядку), на противагу арифметичній теорії, що водночас розглядає \emph{всі цілі елементи}. Тому ясно, що провідник кільця --- особливі точки --- відіграє вирішальну роль. Можна вказати ще дальші аналогії між арифметичною теорією кілець (порядків) і геометричною теорією. Так, геометрична теорія при перетині з ад'юнктними кривими ніде не рахує точки перетину, що припадають на особливі точки, а розглядає лише відмінні від них групи точок. Цьому відповідає те, що Дедекінд\footnote{Über die Anzahl der Idealklassen in den verschiedenen Ordnungen eines endlichen Körpers. (Festschrift Braunschweig 1877); § 3--5. (Теореми, сформульовані в цих параграфах для алгебраїчних чисел, дослівно чинні для алгебраїчних функцій однієї змінної.)} розглядає лише такі ідеали кільця, що є взаємно простими з ідеалом-провідником \(\frf\), і для цієї області ідеалів встановлює всі закони подільності, включно з єдиністю розкладу на прості ідеали.

Також «основна теорема алгебраїчних функцій»\footnote{M. Noether, Über einen Satz aus der Theorie der algebraischen Funktionen. \emph{Math. Ann.} т. 6 (1873).}, що лежить в основі геометричної теорії, або принаймні безпосередній наслідок з неї, у певному сенсі має аналог в арифметичній теорії кілець. Ця теорема дає умови подання
\[
\psi(z,s)=A(z,s)f(z,s)+B(z,s)\varphi(z,s),
\]
де всі величини, що входять, є цілими раціональними за \(s,z\) (загальніше --- цілими й однорідними за \(x_1,x_2,x_3\)). З цих умов випливає,\footnote{Брілль--Нетер, № 55.} що в силу \(f(z,s)=0\) частка \(\frac{\psi}{\varphi}\) завжди дорівнює цілій раціональній функції \(B(z,s)\), якщо вона ад'юнктна до \(f\). Цьому відповідає теорема з п. 3, що лежить в основі всіх цих порівнянь: провідник \(\frf\) кільця, виведеного з \(\delta\), дорівнює ідеалу подвійних точок кривої \(f(z,\delta)=0\). Бо за означенням провідника кожен алгебраїчно цілий елемент, подільний на \(\frf\), належить кільцю, а отже є цілим і раціонально поданим через \(z\) і \(\delta\); тоді як згадана теорема каже, що такий елемент є ад'юнктною функцією.\footnote{Тут, слідом за Дедекіндом, усюди зроблено припущення, що \(\delta\) є цілим алгебраїчним елементом (чого завжди можна досягти перетворенням). Випадок, точніше відповідний геометрії, \(f(z,s)=0\), де \(s\) може також алгебраїчно дробово залежати від \(z\), розглядають Гензель--Ландсберг; пор. Гензель, № 26 і 29.}

\subsection*{7. Порівняння між теоремою про залишок і теоремами теорії ідеалів}
Щойно згаданий наслідок «основної теореми алгебраїчних функцій» у геометричній теорії веде прямо до «теореми про залишок», а отже до основи, на якій будується вся геометрична теорія. У певному сенсі, однак, теорема про залишок знову має елементарніший характер, ніж решта основ, оскільки її можна арифметично сформулювати як прямий наслідок єдиності розкладу ідеалів на прості ідеали, або радше тих фактів, що лежать під цією теоремою розкладу; отже, вона не спирається на вищу теорію кілець. До такого арифметичного формулювання приходять так.

Нехай базову криву \(f(z,\delta)=0\) вибрано так --- чого завжди можна досягти лінійно-дробовим перетворенням змінних, --- що особливі точки, як і скінченно багато груп точок, які розглядаються, лежать у скінченній частині; і нехай, далі, \(\delta\) алгебраїчно ціла над \(z\), чого завжди можна досягти наступним лінійним перетворенням, цілим за змінними. Тоді простому ідеалові \(\frp\), який породжує точку \(P\) абсолютної ріманової поверхні, де \(z=a,\delta=b\), відповідає точка \(z=a,\delta=b\) кривої \(f(z,\delta)=0\); загальніше, ідеал \(\fra=\frp_1^{\,e_1}\cdots\frp_\rho^{\,e_\rho}\) породжує групу точок, що складається з точок \(z=a_i,\delta=b_i\) з відповідною кратністю; і так породжуються всі точки \(f(z,\delta)=0\), що лежать у скінченній частині. Оскільки в точці \(P\) усі функції \(g_1,g_2,\ldots\), подільні на \(\frp\), зникають, і оскільки, навпаки, ця сукупність функцій \(g\) породжує ідеал \(\frp\), відповідна точка на \(f\) є перетином
\[
g_1(z,\delta)=0,\quad g_2(z,\delta)=0,\ldots
\]
з \(f=0\); те саме чинне для групи точок, що відповідає ідеалові \(\fra\). Ба більше, за теоремою, що кожен ідеал є найбільшим спільним дільником двох головних ідеалів, завжди досить двох таких кривих, щоб вирізати групу точок. Зокрема, група точок, відповідна головному ідеалу, вирізається кривою \(g(z,\delta)=0\), і навпаки; головний ідеал відповідає повному перетину.

Дві групи точок \(A\) і \(B\) у геометричній теорії називаються \emph{корезидуальними}, якщо їх можна доповнити однією й тією самою групою точок \(R\) до повного перетину. Якщо ці групи точок відповідно породжуються ідеалами \(\fra,\frb,\frr\), то між цими ідеалами існує відношення
\[
\fra\cdot\frr=(\alpha),\qquad \frb\cdot\frr=(\beta),
\]
яке показує, що ідеали \(\fra\) і \(\frb\) еквівалентні. Навпаки, з еквівалентності двох ідеалів, \(\fra=\eta\frb\), це відношення завжди випливає на підставі теореми, що кожен ідеал можна через множення на інший перетворити на головний ідеал;\footnote{Пор. Дедекінд, Теорія чисел § 181; якщо \(\frr\) вибрано так, що \(\fra\frr\) стає головним ідеалом \((\alpha)\), то \(\frb\frr=\eta^{-1}\cdot(\alpha)\); а оскільки \(\frb\frr\) містить лише цілі елементи, це саме чинне для \(\eta^{-1}\cdot(\alpha)\), який тому також стає головним ідеалом \((\beta)\).} \emph{Отже, еквівалентні ідеали й корезидуальні групи точок є тотожними поняттями.}

Тепер теорема про залишок стверджує: \emph{якщо \(A\) і \(B\) корезидуальні, то їх можна вирізати ад'юнктними кривими з тим самим залишком (поза особливими точками).} Оскільки нулі й полюси функції завжди корезидуальні, цим водночас доведено, що найзагальніші функції поля подаються як частки ад'юнктних форм. Мовою теорії ідеалів теорема про залишок просто зводиться до того, що поряд із \(\fra\) і \(\frb\) також \(\fra\frf\) і \(\frb\frf\) еквівалентні, де \(\frf\) означає ідеал подвійних точок. Бо тоді, за сказаним вище, завжди існує ідеал \(\frm\) такий, що
\[
\fra\frf\frm=(\alpha),\qquad \frb\frf\frm=(\beta)
\]
стають головними ідеалами; \(\alpha=0\) і \(\beta=0\) --- це ад'юнктні криві, які вирізають групи точок \(A\) і \(B\), відповідно їхні підгрупи \(A'\) і \(B'\), відмінні від особливих точок.\footnote{Якщо \(\fra=\frf\fra_1,\ \frb=\frf\frb_1,\ \frr=\frf\frr_1\), де \(\fra_1,\frb_1,\frr_1\) взаємно прості, то вже \(\fra_1\frf\mathfrak n,\ \frb_1\frf\mathfrak n\) стають головними ідеалами, де \(\mathfrak n\) взаємно простий з \(\frf\). Бо кожен ідеал можна перетворити на головний через множення на ідеал, взаємно простий із заданим (Д.--В., с. 214; або Теорія чисел, § 178, теорема XI).} Отже, арифметично доведення теореми про залишок міститься в можливості зіставляти точки простим ідеалам і в доведенні, що еквівалентні ідеали можна множенням на один і той самий ідеал перетворити на головні.

Корезидуальні групи точок не мусять складатися з однакової кількості точок, так само як еквівалентні ідеали не мусять мати однаковий степінь. Наприклад, усі головні ідеали еквівалентні, і всі групи точок, породжені повним перетином, корезидуальні. Ця відмінність від еквівалентних полігонів (дивізорів)\footnote{Полігон \(A\) складається зі скінченної кількості точок абсолютної ріманової поверхні; два полігони \(A\) і \(B\) називаються еквівалентними, якщо вони є нулями і полюсами деякої функції \(\eta\); тоді \(\eta\) допускає подання полігональними частками \(\eta=\frac{A}{B}\). (Д.--В., § 17 і 18.)} пояснюється тим, що точки, в яких \(z\) стає нескінченною, не породжуються простими ідеалами за \(z\). У поданні функції як ідеальної частки \(\eta=\frac{\fra}{\frb}\), що виражає еквівалентність \(\fra\) і \(\frb\), нулі або полюси \(\eta\), де \(z\) стає нескінченною, отже, не з'являються. Так, наприклад, головний ідеал \((\alpha)\) стає добутком простих ідеалів, що відповідають нулям цілої функції \(\alpha\); полюси \(\alpha\) не входять, оскільки \(\alpha\) стає нескінченною лише в точках, де \(z\) стає нескінченною. У цьому сенсі теорія ідеалів є точним аналогом теорії форм у геометричній теорії, де також є лише нулі й немає полюсів.

Перехід від форм до функцій геометрично здійснюється через утворення часток форм однакового порядку; арифметично цьому відповідає перехід від класів ідеалів до інваріантно означених класів полігонів (дивізорів). Бо оскільки тут більше не виділено жодної змінної, подання функції полігональними частками дає всі нулі й полюси. Відношення класу ідеалів (класу корезидуальних груп точок) до класу полігонів таке. Нехай \(\calA\) і \(\calB\) --- полігони, породжені ідеалами \(\fra\) і \(\frb\), а \(\calN\) --- полігон точок, у яких \(z\) стає нескінченною (полігон знаменника \(z\) і цілих функцій від \(z\)). Тоді з еквівалентності \(\fra\) і \(\frb\) випливає еквівалентність \(\calA\) і \(\calB\) тоді й лише тоді, коли \(\fra\) і \(\frb\) мають однаковий степінь; загалом звідси випливає лише еквівалентність \(\calA\calN^\nu\) з \(\calB\calN^{\nu'}\) при \((\nu\ne\nu')\). Отже, поділ ідеалів (відповідно корезидуальних груп точок) на класи є об'єднанням нескінченно багатьох класів полігонів в один клас; його можна також розглядати як поділ полігонів на класи за модулем степеня \(\calN\).\footnote{Гензель--Ландсберг, 25-та лекція, с. 438. --- Якщо подати дві змінні \(z\) і \(s\) полігональними частками, це можна розуміти як бінарне одноріднення кожної з цих змінних; якщо, зокрема, вибрати \(z\) і \(s\) так, щоб вони мали той самий полігон знаменника, виникає звичне в геометрії тернарне одноріднення.}

Теорема про залишок для випадку корезидуальних груп точок, що складаються з різної кількості точок, з'являється лише в суто геометричних питаннях. У застосуванні до алгебраїчних функцій фактично з'являється тільки випадок, відповідний класам полігонів, тобто випадок корезидуальних груп точок однакової кількості. «Повна лінійна система», породжена групою точок, тобто сукупність груп точок, корезидуальних цій групі і з тією самою кількістю точок, точно відповідає класові полігонів, породженому відповідним полігоном; також збігаються поняття «сума двох повних систем» і «добуток двох класів полігонів». Те, що така повна система містить лише скінченно багато лінійно незалежних груп точок, є прямим наслідком теореми про залишок, яка зводить розмірність системи (число лінійно незалежних груп точок) до розмірності всіх ад'юнктних кривих довільного, але фіксованого порядку, що проходять через залишкову групу. Щоб дійти фактичного визначення цієї розмірності, теореми Рімана--Роха, геометрична теорія спершу виводить із теореми про залишок «теорему редукції» (теорему про фіксовані точки);\footnote{Брілль--Нетер, № 58.} а з неї --- теорему Рімана--Роха.

Цілком відповідно арифметична теорія показує, що кожен клас полігонів має скінченну розмірність, і визначає її через спеціальний базис цілих елементів, нормальний базис, означений за допомогою комплементарного модуля; так вона отримує теорему Рімана--Роха як наслідок зв'язку комплементарного модуля з полігоном розгалуження. При цьому у викладі Д.--В. також з'являється теорема редукції,\footnote{Д.--В., § 30, № 2.} хоча й не як основа, як у геометричній теорії. Внаслідок відповідності між базисом та ад'юнктними функціями, згаданої в п. 6, і тут можна встановити певну відповідність допоміжних засобів у двох теоріях. Підсумовуючи, можна сказати, що формування понять у двох теоріях по суті збігаються, хоча, поза формулюванням, вони відрізняються порядком викладу. Обидві теорії виникли незалежно одна від одної; геометрична --- приблизно на десятиліття раніше. Слід ще зауважити, що теорема Рімана--Роха спочатку, у Рімана і Роха, а також у Веєрштраса, з'являється як теорема про ранг: йдеться про ранг матриці \((\varphi_i(x_j))\), де \(\varphi\) пробігають \(p\) ад'юнктних кривих порядку \(n-3\), а \(x\) --- точки залишкової групи.

\subsection*{8. Трансцендентні питання для алгебраїчних функцій і алгебраїчних чисел}
На аналогії між трансцендентними питаннями для алгебраїчних функцій і для алгебраїчних чисел особливо вказував Гільберт;\footnote{Mathematische Probleme, проблема 12. (\emph{Göttinger Nachrichten} 1900).} по суті це питання існування. Так, рімановій постановці питання про існування алгебраїчної функції за заданої ріманової поверхні (отже, також за заданих точок розгалуження) відповідає питання про існування алгебраїчних числових полів із заданою групою і дискримінантом.\footnote{Запропоноване цим трактування алгебраїчних функцій, виходячи з групи, здійснив Р. Кеніґ як спеціальний випадок загальніших питань; пор. особливо: Riemannsche Funktionen- und Differentialsysteme in der Ebene (\emph{J. f. M.} т. 148) і \emph{Math. Ann.} 78, 79. Замість алгебраїчного поля у Кеніґа стоїть спеціальний скінченний модуль щодо \(K(z)\): існує лише «властивість класу».} Однак якщо у випадку алгебраїчних функцій доведення існування можна провести загально, то в області алгебраїчних числових полів воно вдалося лише для спеціальних базових полів (раціональних чисел, квадратичних числових полів) і тільки для абелевих груп.\footnote{Лише простий випадок груп, що з'являються в рівняннях третього і четвертого степенів, а також кілька цілком спеціальних груп у рівняннях 8-го степеня можна розв'язати для довільних базових полів: G. Bucht, Über einige algebraische Körper 8. Grades (\emph{Arkiv f. Math., Astron. och Physik}, т. 6, № 30.) Пор. також E. Noether, Gleichungen mit vorgeschriebener Gruppe і цитовану там дисертацію Зайдельмана. (\emph{Math. Ann.} 78).} Йдеться про теорему Кронекера, що кожне абелеве числове поле над раціональними числами можна скласти з полів коренів з одиниці, і про відповідні теореми Фютера\footnote{\emph{Math. Ann.} 75.} та Гекке\footnote{\emph{Math. Ann.} 76.} для відносно абелевих числових полів. Отже, шукані числові поля фактично даються трансцендентним засобом --- експоненціальною функцією, модулярними функціями, --- на противагу чистому доведенню існування для алгебраїчних функцій.

Доведенню існування алгебраїчних функцій передує існування всюди скінченних інтегралів. Аналогом цього є існування «особливих первинних чисел» алгебраїчного числового поля, означених щодо простого числа \(\ell\), чиї \(\ell\)-ті корені є відносно нерозгалуженими і дають перші будівельні блоки для поля класів.\footnote{Furtwängler, Reziprozitätsgesetze für Potenzreste mit Primzahlexponenten in algebraischen Zahlkörpern. \emph{Math. Ann.} 67 і 72.} Це доведення існування проводиться з істотною допомогою теореми, що в числовому полі завжди існують прості ідеали із заданими лишковими характерами; тому цю теорему можна розглядати як аналог крайової задачі, на якій ґрунтується існування всюди скінченних інтегралів.

Всюди скінченні інтеграли дають трансцендентний критерій для визначення, чи два полігони належать до одного класу (дві групи точок однакової кількості є корезидуальними), а саме через теорему Абеля, яка стверджує, що інтеграли першого роду, взяті вздовж «корезидуальних шляхів», зникають; або через відповідну диференціальну теорему з її оберненням.\footnote{Пор., наприклад, формулювання у M. Noether (\emph{Ann.} 37). Оскільки теорема про залишок історично виросла з теореми Абеля, геометрична теорія говорить про «суму» груп точок і лінійних систем, наслідуючи інтегральні суми; на відміну від «добутку» класів в арифметичній теорії.}

Аналогом теореми Абеля для алгебраїчних числових полів слід вважати «закон взаємності для \(\ell\)-тих степеневих лишків», який для самого поля або принаймні для придатного надполя дає критерій, що два ідеали належать до того самого класу. Його зміст можна сформулювати й так: у цьому надполі особливі первинні числа мають той самий лишковий характер щодо всіх ідеалів того самого класу. Цей останній факт чинний і в самому полі, але там він не збігається із законом взаємності.\footnote{Фуртвенґлер, там само.}

Нарешті, теорія \(p\)-адичних чисел Гензеля дає аналог розкладу алгебраїчних функцій у степеневі ряди; щодо цього див. огляд Гензеля.

\end{document}
